
This work began by observing that different attribution methods give answers that agree on less than 1\% of the identified influential examples. Through systematic validation of four sub-hypotheses, the experiments established that: (1) trade-offs between rank preservation and magnitude fidelity exist in non-convex settings; (2) these trade-offs are absent in convex settings where metrics remain tightly coupled; and (3) the practical consequences are substantial, with methods identifying largely non-overlapping sets of influential examples.

When practitioners ask which attribution method to use, these findings suggest they may first need to ask which quality dimension matters for their application. The search for a universally best attribution method may be complicated in non-convex settings by the structural trade-offs demonstrated here.
